\section{The \instname\ data reduction flow.}

The overall data flow of the \instname\ pipeline is displayed in Figure \ref{fig:reduction_cascade}.

 
\begin{figure}[]
\centering
\includegraphics[width=0.9\textwidth]{../instrument_tutorial/figures/reduction_cascade.jpg}
\caption{The data reduction cascade of the \instname\ workflow.}
\label{fig:reduction_cascade}
\end{figure}

The reduction cascade is organized in tasks, which represent well-defined steps in the process. Tasks can be grouped
inside sub-workflows.
Each task runs a recipe; the detailed description of the algorithms,
input, outputs and recipe parameters used in each recipe are available
in the pipeline manual. Here, we present only the description of most
important features.

The EDPS workflow is designed to execute the tasks that deliver
the final reduced data for each dataset. It can be either the product of a single exposure, or the combination of
multiple exposures. Only calibrations needed by the selected scientific exposures are processed.

It is possible to set EDPS to perform the data reduction until a certain step of the reduction chain (e.g. to reduce
only standard stars, or only flat fields).
This is done by specifying the desired tasks in the field \textbf{Select reduction target} of the \textbf{Raw Data} tab.

The reduction steps of the workflow are listed below. Before starting the reduction,
the parameters of the recipes associated to each task can be configured by pressing the button 
\includegraphics{figures/configure_dataset.jpg} close to each dataset configuration.
See Section \ref{sec:configuration_overview} for more information.

The reduction steps are:

%%%%%%%%%%%%%%%%%%%%%%%%%%%%%%%%%

%%%%%%%%%%%%%%%%%%%%%%%%%%%%%%%%%%%%%%%%%%%%%%%%%%%%%%%%%%%%%%

\subsection{Subworkflow: Dark Background}

Recipes:
\texttt{sph\_ifs\_cal\_background},
\texttt{sph\_ifs\_master\_dark}

This subworkflow generates the instrumental-background and dark-current calibration products. It consists of the following tasks:

\begin{itemize}
    \item \textbf{ifs\_instrument\_background}
    \item \textbf{ifs\_sky\_background}
    \item \textbf{ifs\_dark}
    \item \textbf{ifs\_static\_bad\_pixel\_map}
\end{itemize}

The tasks \textbf{ifs\_instrument\_background} and
\textbf{ifs\_sky\_background} use the recipe
\texttt{sph\_ifs\_cal\_background} to generate instrumental- or
sky-background calibration frames expressed in counts~s$^{-1}$~pixel$^{-1}$.
The output is a multi-extension FITS file containing the background image,
a bad-pixel map, an RMS map, and a weight map.

The task \textbf{ifs\_dark} uses the recipe
\texttt{sph\_ifs\_master\_dark} to generate the detector master dark
(\texttt{IFS\_MASTER\_DARK}) by combining the input raw dark exposures
using the selected collapse algorithm.

The task \textbf{ifs\_static\_bad\_pixel\_map} runs the same recipe,
\texttt{sph\_ifs\_master\_dark}, 
and exports the

\texttt{IFS\_STATIC\_BADPIXELMAP} product, which is also stored in the
second extension of the master dark.

Bad pixels are identified in three successive steps:

\begin{enumerate}
\item Initial thresholding using the
\texttt{min\_acceptable} and
\texttt{max\_acceptable} parameters.
\item Smoothing the master dark with a Gaussian kernel and subtracting the
smoothed frame to remove large-scale variations.
\item Sigma clipping of the residual image using the
\texttt{sigma\_clip} parameter.
\end{enumerate}

The final master dark is written as a multi-extension FITS file containing the dark image, bad-pixel map, RMS map, and a weight map recording the number of contributing raw frames.

When a background correction is required elsewhere in the workflow, EDPS preferentially associates a sky background, then an instrumental background, and finally a master dark if no dedicated background calibration is available.

\paragraph{Customization}

Recipe parameters:

\begin{itemize}
\item \texttt{ifs.master\_dark.coll\_alg}: Collapse algorithm
(\texttt{0} = Mean,
 \texttt{1} = Median,
 \texttt{2} = Clean Mean; default).

\item \texttt{ifs.master\_dark.min\_acceptable},
\texttt{ifs.master\_dark.max\_acceptable}:
Lower and upper pixel-value thresholds used during bad-pixel detection.

\item \texttt{ifs.master\_dark.sigma\_clip}:
Sigma-clipping threshold used to identify bad pixels
(default: \texttt{3}; increasing the value to
\texttt{5} results in fewer rejected pixels).
\end{itemize}

%%%%%%%%%%%%%%%%%%%%%%%%%%%%%%%%%%%%%%%%%%%%%%%%%%%%%%%%%%%%%%

\subsection{Spectra Positions}

Recipe:
\texttt{sph\_ifs\_spectra\_positions}

The task \textbf{ifs\_spectra\_positions} generates the initial mapping
between detector pixels and IFS lenslets, and assigns a first-order
wavelength solution to each detector pixel.

Spectral traces are detected through thresholding and their centroid
positions are measured. These positions are compared with those predicted
by a scaled and shifted lenslet model, allowing refinement of the model
scale and position. Optionally, a two-dimensional distortion model can be
fitted and incorporated into the lenslet model to improve the subsequent
wavelength calibration.

The main output product is \texttt{IFS\_SPECPOS}, a pixel description table (PDT), stored as a multi-extension FITS file containing:

\begin{itemize}
\item Wavelength
\item Spectral region identifier
\item Lenslet identifier
\item Wavelength width
\item Second derivative of the wavelength
\item Illumination fraction
\end{itemize}

%\paragraph{Customization}

%Recipe parameters:

%\begin{itemize}
%\item \texttt{ifs.spectra\_positions.coll\_alg}:
%Collapse algorithm
%(\texttt{0} = Mean,
% \texttt{1} = Median,
% \texttt{2} = Clean Mean; default).
%\end{itemize}
% Parameter not found in EDPS code

%%%%%%%%%%%%%%%%%%%%%%%%%%%%%%%%%%%%%%%%%%%%%%%%%%%%%%%%%%%%%%

\subsection{Wavelength Calibration}

Recipe:
\texttt{sph\_ifs\_wave\_calib}

The task \textbf{ifs\_wavelength\_calibration} performs the wavelength
calibration by refining the pixel-to-wavelength associations produced by
\textbf{ifs\_spectra\_positions} using wavelength calibration exposures
and the existing lenslet model.

One-dimensional spectra are extracted for each spectral region and the
positions of known calibration lines are measured using a flux-weighted
centroid within a configurable fitting window. The measured line positions
are fitted with a polynomial defining the wavelength solution for each
spectrum. If the derived solution is inconsistent with the expected
dispersion, the original model wavelengths are retained or the affected
spectral region is flagged as invalid.

From the final wavelength solutions, the recipe computes quality-control
parameters for each spectrum, including the minimum, maximum, and central
wavelengths, as well as the spectral resolving power derived from the local wavelength gradient.

The main output product is \texttt{IFS\_WAVECALIB}, which contains the
updated wavelength calibration derived from the initial
\texttt{IFS\_SPECPOS} product.

%\paragraph{Customization}

%Recipe parameters:

%\begin{itemize}
%\item \texttt{ifs.wave\_calib.coll\_alg}:
%Collapse algorithm
%(\texttt{0} = Mean,
% \texttt{1} = Median,
% \texttt{2} = Clean Mean; default).
%\end{itemize}

%%%%%%%%%%%%%%%%%%%%%%%%%%%%%%%%%%%%%%%%%%%%%%%%%%%%%%%%%%%%%%

\subsection{Subworkflow: Detector Flat}

Recipe:
\texttt{sph\_ifs\_master\_detector\_flat}

This subworkflow generates detector flat-field calibrations from narrow-
and broadband lamp exposures. It consists of the following tasks:

\begin{itemize}
\item \textbf{ifs\_det\_flat\_narrow\_band1}
\item \textbf{ifs\_det\_flat\_narrow\_band2}
\item \textbf{ifs\_det\_flat\_narrow\_band3}
\item \textbf{ifs\_det\_flat\_narrow\_band4} (only associated for the \texttt{PRI\_YJH} disperser.)
\item \textbf{ifs\_det\_flat\_broad\_band}
\end{itemize}

The recipe produces several detector flat-field components, including a
standard detector flat, a large-scale (smoothed) flat, and a
preamplifier-correction flat. In most applications, the standard
\texttt{IFS\_MASTER\_DFF\_LONG*} products are sufficient to correct
pixel-to-pixel detector response variations.

Raw frames are first corrected for dark current and preamplifier
variations. The detector flat is then derived by fitting, for each pixel,
a linear relation between its signal and the mean illumination level of
the corresponding exposure. The slope of this fit represents the relative
pixel response.

Optional products include a detector non-linearity map and a smoothed
large-scale flat-field image.

\paragraph{Customization}

Recipe parameters:

\begin{itemize}
%\item \texttt{ifs.master\_detector\_flat.coll\_alg}:
%Collapse algorithm
%(\texttt{0} = Mean,
% \texttt{1} = Median,
% \texttt{2} = Clean Mean; default).

\item \texttt{ifs.master\_detector\_flat.robust\_fit}:
Controls the fitting method. By default
(\texttt{TRUE}) a maximum-likelihood fit is used. Setting the parameter to
\texttt{FALSE} enables a robust linear fit that is more resistant to
outliers such as cosmic rays.
\end{itemize}

%%%%%%%%%%%%%%%%%%%%%%%%%%%%%%%%%%%%%%%%%%%%%%%%%%%%%%%%%%%%%%

\subsection{IFU Flat}

Recipe:
\texttt{sph\_ifs\_instrument\_flat}

The task \textbf{ifs\_ifu\_flat} uses the pipeline recipe to generate
flat-field calibrations for the IFS. The underlying recipe supports two operating modes. Within the EDPS workflow, the IFU-flat mode is used to generate the calibration products required by subsequent reductions.

In detector-flat (or total-flat) mode, raw calibration frames are combined
to generate a flat field that includes the detector response. This product
can subsequently be used by the
\textbf{ifs\_spectra\_positions} and
\textbf{ifs\_wavelength\_calibration} tasks.

In IFU-flat mode, the recipe removes the detector response using the pixel
description table (PDT) and wavelength calibration products in order to
isolate the transmission of the IFU lenslet array.

A wavelength-dependent ``super-flat'' is first constructed from the master
detector flats and applied to the combined calibration frames. Using the
lenslet model provided by the wavelength calibration, spectra are extracted
for every lenslet and collapsed along the wavelength direction to derive a
single flat-field value for each lenslet.

The main outputs are:

\begin{itemize}
\item \texttt{IFS\_INSTRUMENT\_FLAT\_FIELD}: detector flat-field product
containing the flat-field image, bad-pixel map, RMS map, and weight map.

\item \texttt{IFS\_IFU\_FLAT\_FIELD}: IFU flat-field product containing
the flat-field image, bad-pixel map, RMS map, weight map, together with
a binary-table extension listing the lenslet flat-field values.

\end{itemize}

%\paragraph{Customization}

%Recipe parameters:

%\begin{itemize}
%\item \texttt{ifs.instrument\_flat.coll\_alg}:
%Collapse algorithm
%(\texttt{0} = Mean,
% \texttt{1} = Median,
% \texttt{2} = Clean Mean; default).
%\end{itemize}

%%%%%%%%%%%%%%%%%%%%%%%%%%%%%%%%%%%%%%%%%%%%%%%%%%%%%%%%%%%%%%

\subsection{Distortion Map}

Recipe:
\texttt{sph\_ifs\_distortion\_map}

The task \textbf{ifs\_distortion\_map} measures the geometric distortion
of the SPHERE IFS lenslet grid.

Raw calibration frames are first reduced in a manner similar to science
data using the master dark. Instrumental background, wavelength calibration, detector flats and IFU flat may additionally be used when available. The resulting monochromatic images are collapsed along the
wavelength axis and point sources are identified using a user-defined
detection threshold.

If no reference point pattern is provided, the recipe derives one from the
detected sources and writes it as an output product. Otherwise, the supplied
reference pattern is used directly.

The distortion is measured by comparing the detected and expected source
positions. Vectors larger than a configurable threshold are rejected before
the remaining offsets are fitted with a two-dimensional polynomial model.

Optionally, the derived distortion solution can be applied to the data to
produce a residual distortion map, providing a diagnostic of the quality of
the distortion correction.

The outputs are:

\begin{itemize}
\item \texttt{IFS\_POINT\_PATTERN}: reference point pattern that may be
reused as input for subsequent executions of the recipe.

\item \texttt{IFS\_DISTORTION\_MAP}: multi-extension FITS file containing
the distortion maps in both detector coordinates. For each coordinate, the
product includes the distortion image, bad-pixel map, RMS map, and weight
map.
\end{itemize}

\paragraph{Customization}

Recipe parameters:

\begin{itemize}
%\item \texttt{ifs.distortion\_map.coll\_alg}:
%Collapse algorithm
%(\texttt{0} = Mean,
% \texttt{1} = Median,
% \texttt{2} = Clean Mean; default).

\item \texttt{ifs.distortion\_map.threshold}:
Detection threshold used to identify calibration spots
(default: \texttt{3.0}).
\end{itemize}

%%%%%%%%%%%%%%%%%%%%%%%%%%%%%%%%%%%%%%%%%%%%%%%%%%%%%%%%%%%%%%

\subsection{Subworkflow: Science Data Reduction}

Recipe:
\texttt{sph\_ifs\_science\_dr}

This subworkflow performs the reduction of IFS science observations. It
consists of the following tasks:

\begin{itemize}
\item \textbf{ifs\_science\_flux}
\item \textbf{ifs\_coronagraph\_center}
\item \textbf{ifs\_science}
\end{itemize}

The workflow reduces science observations with sky-, instrument background-, or dark-subtraction and preamplifier stripe correction.

It also processes \textbf{FLUX} frames obtained during coronagraphic
observations, in which the telescope is offset to move the target away
from the coronagraphic mask in order to measure the stellar flux.

The workflow also processes \textbf{CENTER} frames used to determine the stellar position behind the coronagraph. The resulting products are primarily intended for later analysis rather than subsequent EDPS processing.

Large-scale flat-field variations are corrected using a wavelength-dependent ``super-flat'' derived from multi-colour lamp flats and the wavelength calibration, while small-scale detector response variations are corrected using the broadband detector flat.

The recipe also supports automatic post-processing through spectral
deconvolution (Mesa et al. 2015) and Angular Differential Imaging (ADI),
which can be enabled or disabled through recipe parameters.

The processing is currently done file-by-file because the pipeline does not stack science data.

\paragraph{Customization}

Workflow parameters:

\begin{itemize}
\item \texttt{force\_distortion\_correction}: Force the association of distortion calibration to IFU scientific task. Possible values are \texttt{TRUE}/\texttt{FALSE}. QC0 process requires this parameter set to \texttt{TRUE}. In the current EDPS workflow, the associated distortion calibration is not propagated to the science recipe.
\end{itemize}

Recipe parameters:

\begin{itemize}
%\item \texttt{ifs.science\_dr.coll\_alg}:
%Collapse algorithm
%(\texttt{0} = Mean,
% \texttt{1} = Median,
% \texttt{2} = Clean Mean; default).

\item \texttt{ifs.science\_dr.use\_adi}:
Controls the use of Angular Differential Imaging
(\texttt{0} = disabled,
 \texttt{1} = always enabled,
 \texttt{2} = automatically enabled when the total field rotation exceeds
 
 \texttt{ifs.science\_dr.min\_adi\_angle}).

\item \texttt{ifs.science\_dr.min\_adi\_angle}:
Minimum field rotation required for the automatic ADI mode.

\item \texttt{ifs.science\_dr.spec\_deconv}:
Controls the use of spectral deconvolution
(\texttt{0} = disabled,
 \texttt{1} = enabled).
\end{itemize}

%%%%%%%%%%%%%%%%%%%%%%%%%%%%%%%%%%%%%%%%%%%%%%%%%%%%%%%%%%%%%%

\subsection{Subworkflow: Standard Data Reduction}

Recipe:
\texttt{sph\_ifs\_science\_dr}

This subworkflow consists of the following tasks:

\begin{itemize}
\item \textbf{ifs\_standard\_astrometry}
\item \textbf{ifs\_standard\_flux}
\end{itemize}

At present, astrometric-standard and flux-standard observations are reduced using the same workflow and recipe as science observations. Consequently, the same calibration steps, reduction algorithms, and recipe parameters described in the previous section are applied.
